\chapter{Experimental Evaluation}

\section{Evaluation Dataset}
The evaluation uses a set of realistic student-style queries about UniTN bureaucracy. Queries should cover different categories, such as study plans, graduation, internships, Erasmus, taxes, exam procedures, and course-specific rules.

Each query is evaluated in two conditions:

\begin{enumerate}
    \item BM25 search using the original query;
    \item BM25 search using the LLM-rewritten query.
\end{enumerate}

\section{Manual Relevance Judgement}
Retrieved results are manually labelled according to their relevance to the query. A simple three-level scale is used.

\begin{table}[h]
\centering
\begin{tabular}{clp{0.58\linewidth}}
\toprule
\textbf{Label} & \textbf{Name} & \textbf{Meaning} \\
\midrule
2 & Relevant & The result is an official source that directly answers the query. \\
1 & Partially relevant & The result is related but incomplete, ambiguous, or only indirectly useful. \\
0 & Not relevant & The result does not answer the query or is unrelated. \\
\bottomrule
\end{tabular}
\caption{Manual relevance scale.}
\label{tab:relevance-scale}
\end{table}

The primary metrics treat label 2 as relevant. A secondary analysis can also report results when labels 1 and 2 are both treated as useful.

\section{Metrics}

\subsection{Success@1}
Success@1 measures whether the first retrieved result is relevant:

\begin{equation}
Success@1(q) =
\begin{cases}
1, & \text{if the result at rank 1 is relevant} \\
0, & \text{otherwise.}
\end{cases}
\end{equation}

This metric is useful because users often inspect the first result first.

\subsection{Success@5}
Success@5 measures whether at least one relevant result appears in the first five results:

\begin{equation}
Success@5(q) =
\begin{cases}
1, & \text{if any result at ranks 1--5 is relevant} \\
0, & \text{otherwise.}
\end{cases}
\end{equation}

This metric is useful because the interface can show multiple candidates and the system may still be useful if the correct source appears within the first page of results.

\subsection{Mean Reciprocal Rank}
Mean Reciprocal Rank (MRR) rewards systems that place the first relevant result higher in the ranking:

\begin{equation}
MRR = \frac{1}{|Q|} \sum_{q \in Q} \frac{1}{rank_q}
\end{equation}

where $rank_q$ is the rank position of the first relevant result for query $q$. If no relevant result is found, the reciprocal rank is defined as 0.

\section{Experimental Comparison}
The main comparison is summarized in Table~\ref{tab:experimental-conditions}.

\begin{table}[h]
\centering
\begin{tabular}{lll}
\toprule
\textbf{Condition} & \textbf{Input to BM25} & \textbf{Purpose} \\
\midrule
Baseline & Original query & Measure raw BM25 performance \\
Rewrite & LLM-rewritten query & Measure effect of semantic rewriting \\
\bottomrule
\end{tabular}
\caption{Experimental conditions.}
\label{tab:experimental-conditions}
\end{table}

\section{Planned Result Tables}
The following table will be filled after running the evaluation.

\begin{table}[h]
\centering
\begin{tabular}{lccc}
\toprule
\textbf{Condition} & \textbf{Success@1} & \textbf{Success@5} & \textbf{MRR} \\
\midrule
BM25 original query & TODO & TODO & TODO \\
BM25 rewritten query & TODO & TODO & TODO \\
\bottomrule
\end{tabular}
\caption{Retrieval performance comparison.}
\label{tab:results-main}
\end{table}
\par
